% S11. Réplicas por región: leader-follower, lag acotado, gaps, checksum, conflation.
\begin{tikzpicture}[x=1cm,y=1cm]
  \path[use as bounding box] (0,0.05) rectangle (13.5,4.6);
  \tikzset{rl/.style={font=\scriptsize,text width=1.2cm,inner sep=2pt}}
  % ---------------- Diagrama ----------------
  \node[acento,font=\scriptsize,text width=3.6cm] (or) at (3.3,4.05) {\textbf{Origen = líder}\\numera cada mensaje};
  \draw[colNeutro!60,rounded corners=3pt] (0.05,0.1) rectangle (3.2,2.95);
  \draw[colNeutro!60,rounded corners=3pt] (3.45,0.1) rectangle (6.6,2.95);
  \node[font=\scriptsize\bfseries] at (1.62,2.72) {Región A};
  \node[font=\scriptsize\bfseries] at (5.02,2.72) {Región B};
  \node[bueno,rl] (a1) at (0.85,1.85) {Relay 1\\sano};
  \node[malo,rl]  (a2) at (2.4,1.85) {Relay 2\\lag $>$ 500\,ms};
  \node[bueno,rl] (b1) at (4.25,1.85) {Relay 3\\sano};
  \node[bueno,rl] (b2) at (5.8,1.85) {Relay 4\\sano};
  \foreach \r in {a1,a2,b1,b2} \draw[flow,colAcento] (or.south) -- (\r.north);
  \node[font=\tiny,text=colAcento,fill=white,inner sep=1pt] at (3.3,3.25) {feed = log de replicación};
  \node[cliente] (cl) at (2.4,0.55) {cliente};
  \draw[colWake,dashed] (a2.south) -- (cl.north);
  \draw[rec,colReq] (cl.west) to[out=180,in=270] node[pos=0.35,below,font=\tiny,text=colReq] {reanuda aquí} (a1.south);
  \foreach \x in {4.25,5.8} {
    \node[cliente,minimum width=0.8cm] at (\x,0.55) {cl.};
    \draw[flow,colAcento,thin] (\x,1.3) -- (\x,0.83);
  }
  % ---------------- Mecanismos ----------------
  \tikzset{mec/.style={box,font=\scriptsize,text width=6.2cm,align=left,anchor=west}}
  \node[mec] at (6.95,4.15) {\textbf{1. Un solo orden:} solo el origen numera; todos aplican lo mismo.};
  \node[mec] at (6.95,3.25) {\textbf{2. Lag acotado:} si un relay se atrasa $>$\,500\,ms, sus clientes van a otro relay.};
  \node[mec] at (6.95,2.35) {\textbf{3. Gaps:} redirige o espera el snapshot del feed (cada 1\,s); nunca pide al origen.};
  \node[mec] at (6.95,1.45) {\textbf{4. Checksum} de los 10 mejores niveles: detecta copias que divergen.};
  \node[mec,bueno,align=left] at (6.95,0.55) {\textbf{5. Conflation:} al cliente lento no se le corta; recibe solo el último valor.};
\end{tikzpicture}
